% ============================================================
\chapter{Espaces Bitopologiques}
\label{ch:bitopologiques}
% ============================================================

En 1963, John C.\ Kelly propose une id\'ee audacieuse~: munir un m\^eme ensemble de \emph{deux} topologies distinctes et \'etudier leur interaction. L'id\'ee n'est pas gratuite --- elle \'emerge naturellement des quasi-m\'etriques, o\`u la topologie \og avant \fg et la topologie \og arri\`ere \fg cohabitent, et de la th\'eorie de l'ordre, o\`u les topologies haute et basse sur un ensemble ordonn\'e capturent des aspects compl\'ementaires. La dualit\'e de Nachbin entre espaces ordonn\'es compacts et espaces stablement compacts montre que cette structure bitopologique n'est pas un artefact, mais un ph\'enom\`ene profond.

\begin{intuition}
Un espace bitopologique porte \emph{deux} topologies sur le m\^eme ensemble, capturant
une dualit\'e naturelle. Cette structure appara\^it naturellement via les quasi-m\'etriques
(topologie directe et conjugu\'ee), en th\'eorie de l'ordre (topologies haute et basse),
et dans la dualit\'e de Nachbin entre espaces ordonn\'es compacts et espaces stablement
compacts.
\end{intuition}

% ------------------------------------------------------------
\section{D\'efinitions fondamentales}
% ------------------------------------------------------------

\begin{definition}[Espace bitopologique]
Un \emph{espace bitopologique} est un triplet $(X, \tau_1, \tau_2)$ o\`u $X$ est un
ensemble et $\tau_1, \tau_2$ sont deux topologies sur $X$. On appelle $\tau_1$ la topologie
\emph{primaire} et $\tau_2$ la topologie \emph{secondaire} (ou \emph{duale}).
\end{definition}

\begin{definition}[Application bicontinue]
Soient $(X, \tau_1, \tau_2)$ et $(Y, \sigma_1, \sigma_2)$ deux espaces bitopologiques. Une
application $f\colon X \to Y$ est \emph{bicontinue} si $f$ est continue de $(X,\tau_1)$
vers $(Y,\sigma_1)$ et de $(X,\tau_2)$ vers $(Y,\sigma_2)$.
\end{definition}

\begin{definition}[Cat\'egorie $\mathbf{BiTop}$]
La cat\'egorie $\mathbf{BiTop}$ a pour objets les espaces bitopologiques et pour
morphismes les applications bicontinues. Un \emph{bi-hom\'eomorphisme} est un isomorphisme
dans $\mathbf{BiTop}$.
\end{definition}

% ------------------------------------------------------------
\section{Exemples fondamentaux}
% ------------------------------------------------------------

\begin{exemple}[Espace bitopologique d'une quasi-m\'etrique]
Soit $(X,d)$ un espace quasi-m\'etrique. L'espace bitopologique associ\'e est
$(X, \tau_d, \tau_{d^{-1}})$, o\`u $\tau_d$ est la topologie directe et $\tau_{d^{-1}}$ la
topologie de la conjugu\'ee.
\end{exemple}

\begin{exemple}[Espaces ordonn\'es]
Soit $(X,\leq)$ un ensemble ordonn\'e. L'espace bitopologique d'Alexandrov est
$(X, \tau_{\mathrm{up}}, \tau_{\mathrm{down}})$, o\`u $\tau_{\mathrm{up}}$ est la topologie
d'Alexandrov (ouverts = parties hautes) et $\tau_{\mathrm{down}}$ la topologie duale
(ouverts = parties basses).
\end{exemple}

\begin{exemple}[Droite r\'eelle ordonn\'ee]
Sur $\R$, la topologie des ouverts \`a droite $\{(a,+\infty) : a \in \R\} \cup
\{\emptyset, \R\}$ et la topologie des ouverts \`a gauche
$\{(-\infty, b) : b \in \R\} \cup \{\emptyset, \R\}$ d\'efinissent un espace
bitopologique. Leur supremum est la topologie usuelle de $\R$.
\end{exemple}

% ------------------------------------------------------------
\section{Axiomes de s\'eparation bitopologiques}
% ------------------------------------------------------------

\begin{definition}[S\'eparation de Kelly]
Un espace bitopologique $(X, \tau_1, \tau_2)$ est~:
\begin{itemize}
  \item \emph{Pairwise $T_0$} si pour tous $x \neq y$, il existe un ouvert de $\tau_1$ ou
    de $\tau_2$ contenant l'un mais pas l'autre.
  \item \emph{Pairwise $T_1$} si pour tous $x \neq y$, il existe un $\tau_1$-ouvert
    contenant $x$ mais pas $y$, et un $\tau_2$-ouvert contenant $y$ mais pas $x$ (ou
    inversement).
  \item \emph{Pairwise $T_2$} (ou \emph{pairwise Hausdorff}) si pour tous $x \neq y$, il
    existe un $\tau_1$-ouvert $U \ni x$ et un $\tau_2$-ouvert $V \ni y$ avec
    $U \cap V = \emptyset$.
\end{itemize}
\end{definition}

\begin{theoreme}[Pairwise $T_2$ et quasi-m\'etriques]
Si $(X,d)$ est une quasi-m\'etrique $T_0$-s\'epar\'ee, alors $(X,\tau_d, \tau_{d^{-1}})$
est pairwise $T_2$.
\end{theoreme}

\begin{proof}
Soient $x \neq y$. Alors $d(x,y) > 0$ ou $d(y,x) > 0$. Supposons $d(x,y) > 0$ et posons
$\varepsilon = d(x,y)/2 > 0$. Alors $U = B_d(x,\varepsilon)$ est un $\tau_d$-ouvert
contenant $x$ et $V = B_{d^{-1}}(y,\varepsilon)$ est un $\tau_{d^{-1}}$-ouvert contenant
$y$. Si $z \in U \cap V$, alors $d(x,z) < \varepsilon$ et $d^{-1}(y,z) = d(z,y) <
\varepsilon$, donc $d(x,y) \leq d(x,z) + d(z,y) < 2\varepsilon = d(x,y)$, contradiction.
Donc $U \cap V = \emptyset$.
\end{proof}

% ------------------------------------------------------------
\section{Normalit\'e bitopologique}
% ------------------------------------------------------------

\begin{definition}[Normalit\'e pairwise]
Un espace bitopologique $(X,\tau_1,\tau_2)$ est \emph{pairwise normal} si pour tout
$\tau_1$-ferm\'e $F_1$ et tout $\tau_2$-ferm\'e $F_2$ avec $F_1 \cap F_2 = \emptyset$,
il existe un $\tau_2$-ouvert $U_2 \supseteq F_1$ et un $\tau_1$-ouvert $U_1 \supseteq F_2$
avec $U_1 \cap U_2 = \emptyset$.
\end{definition}

\begin{theoreme}[Lemme d'Urysohn bitopologique]
Si $(X,\tau_1,\tau_2)$ est pairwise normal et pairwise $T_1$, alors pour tout
$\tau_1$-ferm\'e $F_1$ et tout $\tau_2$-ferm\'e $F_2$ disjoints, il existe une fonction
$f\colon X \to [0,1]$ telle que~:
\begin{itemize}
  \item $f \equiv 0$ sur $F_1$ et $f \equiv 1$ sur $F_2$;
  \item $f$ est semi-continue inf\'erieurement pour $\tau_1$ et semi-continue
    sup\'erieurement pour $\tau_2$.
\end{itemize}
\end{theoreme}

\begin{proof}
La preuve suit le sch\'ema classique du lemme d'Urysohn. On construit par r\'ecurrence
sur les dyadiques $r \in [0,1] \cap \Z[1/2]$ une famille d'ensembles $(A_r)$ tels que~:
\begin{itemize}
  \item $F_1 \subseteq A_0$ et $A_1 = X \setminus F_2$;
  \item $A_r$ est $\tau_2$-ouvert pour $r < 1$;
  \item $r < s \Rightarrow \overline{A_r}^{\tau_1} \subseteq A_s$.
\end{itemize}
L'existence de ces ensembles d\'ecoule de la normalit\'e pairwise. On d\'efinit
$f(x) = \inf\{r : x \in A_r\}$ et on v\'erifie les propri\'et\'es de semi-continuit\'e.
\end{proof}

% ------------------------------------------------------------
\section{Compacit\'e bitopologique}
% ------------------------------------------------------------

\begin{definition}[Compacit\'e jointe]
Un espace bitopologique $(X,\tau_1,\tau_2)$ est \emph{joinly compact} (ou
\emph{conjointement compact}) si l'espace $(X, \tau_1 \vee \tau_2)$ est compact.
\end{definition}

\begin{definition}[Compacit\'e pairwise]
Un espace bitopologique est \emph{pairwise compact} si tout recouvrement form\'e d'ouverts
de $\tau_1$ et de $\tau_2$ admet un sous-recouvrement fini.
\end{definition}

\begin{proposition}
La compacit\'e pairwise implique la compacit\'e jointe. La r\'eciproque est vraie si
$\tau_1$ et $\tau_2$ forment une <<~paire compatible~>>.
\end{proposition}

% ------------------------------------------------------------
\section{La dualit\'e de Nachbin}
% ------------------------------------------------------------

\begin{definition}[Espace ordonn\'e compact]
Un \emph{espace ordonn\'e compact} (au sens de Nachbin) est un triplet $(X, \tau, \leq)$
o\`u~:
\begin{enumerate}
  \item $(X,\tau)$ est un espace compact de Hausdorff;
  \item $\leq$ est un ordre partiel ferm\'e dans $X \times X$ (pour la topologie produit).
\end{enumerate}
\end{definition}

\begin{theoreme}[Dualit\'e de Nachbin]
Il y a une \'equivalence de cat\'egories entre~:
\begin{enumerate}
  \item la cat\'egorie des espaces ordonn\'es compacts (avec applications continues
    croissantes);
  \item la cat\'egorie des espaces bitopologiques \emph{pairwise compact pairwise $T_2$
    joinly $T_1$} (avec applications bicontinues).
\end{enumerate}
\end{theoreme}

\begin{proof}[Id\'ee de la preuve]
\noindent$(\Rightarrow)$ Soit $(X,\tau,\leq)$ un espace ordonn\'e compact. On d\'efinit~:
\begin{itemize}
  \item $\tau_1 = $ la topologie des ouverts hauts (topologie sup\'erieure associ\'ee
    \`a $\leq$ dans $\tau$);
  \item $\tau_2 = $ la topologie des ouverts bas.
\end{itemize}
Plus pr\'ecis\'ement, $\tau_1$ est engendr\'ee par les ouverts $U \in \tau$ tels que
$U = \uparrow U$ et $\tau_2$ par les ouverts $V \in \tau$ tels que $V = \downarrow V$.
On a $\tau_1 \vee \tau_2 = \tau$ par le th\'eor\`eme de Nachbin.

$(\Leftarrow)$ \'Etant donn\'e un espace bitopologique $(X,\tau_1,\tau_2)$ de la
cat\'egorie d\'ecrite, on pose $\tau = \tau_1 \vee \tau_2$ et l'ordre est
$x \leq y \Leftrightarrow x \in \overline{\{y\}}^{\tau_1}$.
\end{proof}

% ------------------------------------------------------------
\section{Lien avec les espaces stablement compacts}
% ------------------------------------------------------------

\begin{theoreme}[Espaces stablement compacts et bitopologiques]
Soit $X$ un espace stablement compact. Alors~:
\begin{enumerate}
  \item La topologie co-compacte $\tau^d$ (engendr\'ee par les compl\'ementaires des compacts
    satur\'es) d\'etermine un espace bitopologique $(X, \tau, \tau^d)$.
  \item Le triplet $(X, \tau \vee \tau^d, \leq)$ est un espace ordonn\'e compact.
  \item Cette construction \'etablit une \'equivalence entre la cat\'egorie des espaces
    stablement compacts et une sous-cat\'egorie des espaces ordonn\'es compacts.
\end{enumerate}
\end{theoreme}

\begin{remarque}
C'est l'un des r\'esultats les plus profonds de la topologie asym\'etrique~: les espaces
stablement compacts, les espaces ordonn\'es compacts et certains espaces bitopologiques
sont trois facettes d'un m\^eme ph\'enom\`ene.
\end{remarque}

% ------------------------------------------------------------
\section{Foncteurs entre $\mathbf{BiTop}$ et $\mathbf{Top}$}
% ------------------------------------------------------------

\begin{proposition}[Foncteurs d'oubli et de sym\'etrisation]
On dispose des foncteurs suivants~:
\begin{enumerate}
  \item $U_1, U_2\colon \mathbf{BiTop} \to \mathbf{Top}$, projection sur chaque topologie.
  \item $J\colon \mathbf{BiTop} \to \mathbf{Top}$, envoyant $(X,\tau_1,\tau_2)$ sur
    $(X, \tau_1 \vee \tau_2)$ (topologie jointe).
  \item $\Delta\colon \mathbf{Top} \to \mathbf{BiTop}$, envoyant $(X,\tau)$ sur
    $(X,\tau,\tau)$ (doublement diagonal).
\end{enumerate}
Le foncteur $\Delta$ est adjoint \`a gauche et \`a droite de $J$ dans certaines
sous-cat\'egories.
\end{proposition}

% ------------------------------------------------------------
\section{Espaces bitopologiques et treillis}
% ------------------------------------------------------------

\begin{definition}[Biframe]
Un \emph{biframe} est un triplet $(L, L_1, L_2)$ o\`u $L$ est un cadre et $L_1, L_2$
sont des sous-cadres de $L$ qui engendrent $L$~: chaque \'el\'ement de $L$ est un sup
d'infs finis d'\'el\'ements de $L_1 \cup L_2$.
\end{definition}

\begin{remarque}
Les biframes sont l'analogue <<~sans points~>> des espaces bitopologiques, de m\^eme que
les cadres sont l'analogue des espaces topologiques. La th\'eorie des bilocales
d\'evelopp\'ee par Banaschewski, Br\"ummer et Hardie fournit un cadre alg\'ebrique riche
pour la bitopologie.
\end{remarque}

% ------------------------------------------------------------
\section{Exercices}
% ------------------------------------------------------------

\begin{exercice}
Montrer que pour toute quasi-m\'etrique $d$, l'espace bitopologique
$(X,\tau_d,\tau_{d^{-1}})$ est pairwise $T_0$ si et seulement si $d$ est $T_0$-s\'epar\'ee.
\end{exercice}

\begin{exercice}
Soit $(X,\tau_1,\tau_2)$ un espace bitopologique pairwise $T_2$. Montrer que la topologie
jointe $\tau_1 \vee \tau_2$ est de Hausdorff.
\end{exercice}

\begin{exercice}
Montrer que tout espace ordonn\'e compact est pairwise normal.
\end{exercice}

\begin{exercice}
Soit $X = [0,1]$ avec l'ordre usuel. D\'ecrire explicitement les topologies
$\tau_{\mathrm{up}}$, $\tau_{\mathrm{down}}$, la topologie jointe, et v\'erifier que
$(X, \tau_{\mathrm{up}} \vee \tau_{\mathrm{down}}, \leq)$ est un espace ordonn\'e compact.
\end{exercice}

\begin{exercice}
Montrer que la cat\'egorie $\mathbf{BiTop}$ est compl\`ete et cocompl\`ete. D\'ecrire les
produits et coproduits.
\end{exercice}

\begin{exercice}[Dualit\'e de Priestley]
Soit $D$ un treillis distributif born\'e et $X$ l'espace de Priestley associ\'e (espace
ordonn\'e compact totalement s\'epar\'e par l'ordre). D\'ecrire l'espace bitopologique
correspondant via la dualit\'e de Nachbin et montrer qu'on retrouve l'espace spectral de
la dualit\'e de Stone.
\end{exercice}
